% TEMPLATE-MILC-v3.tex - plantilla maestra para todas las guias MILC v3
% Ver scripts/generadores/build-guia-milc.py para la lista completa de
% claves esperadas. El script valida que no queden placeholders sin
% reemplazar antes de escribir el archivo final.

\documentclass[11pt,letterpaper]{article}
\usepackage[margin=1.75cm]{geometry}
\usepackage[table]{xcolor}
\usepackage{fontspec}
\usepackage[spanish]{babel}
\usepackage{tikz}
\usepackage[most]{tcolorbox}
\usepackage{tabularx}
\usepackage{array}
\usepackage{fancyhdr}
\usepackage{titlesec}
\usepackage{ragged2e}
\usepackage{enumitem}
\usepackage{microtype}

\setmainfont{Helvetica Neue}
\setsansfont{Helvetica Neue}
\setlength{\parindent}{0pt}
\setlength{\parskip}{6pt}
\hyphenpenalty=200
\exhyphenpenalty=200
\pretolerance=400
\tolerance=2000
\setlength{\emergencystretch}{1em}
\renewcommand{\arraystretch}{1.25}
\setlist{leftmargin=5mm,itemsep=1mm,topsep=1mm}
\newcolumntype{Y}{>{\RaggedRight\arraybackslash}X}

\definecolor{milcMagenta}{HTML}{E6007E}
\definecolor{milcVino}{HTML}{5A0038}
\definecolor{milcTurquesa}{HTML}{0093A5}
\definecolor{milcVerde}{HTML}{58A923}
\definecolor{milcMostaza}{HTML}{E5B400}
\definecolor{milcOcre}{HTML}{B86600}
\definecolor{milcNegro}{HTML}{241816}
\definecolor{milcGris}{HTML}{F7F7F4}
\definecolor{milcLinea}{HTML}{E8E2DD}
\definecolor{milcGradePrimary}{HTML}{0093A5}
\definecolor{milcGradeDark}{HTML}{00525C}
\definecolor{milcGradeSoft}{HTML}{E0F3F5}
\definecolor{milcGradeAccent}{HTML}{CFE7FF}
\definecolor{milcGradeText}{HTML}{FFFFFF}
\color{milcNegro}

\pagestyle{fancy}
\fancyhf{}
\lhead{\small Tecnología e Informática · Grado 8}
\rhead{\small Guía 11 · MILC v3}
\cfoot{\small\thepage}
\renewcommand{\headrulewidth}{0pt}

\newtcolorbox{titlebox}[1]{
  enhanced,colback=#1,colframe=#1,arc=7mm,boxrule=0pt,
  left=7mm,right=7mm,top=3mm,bottom=3mm,
  before skip=8pt,after skip=8pt,
  fontupper=\sffamily\bfseries\Large\color{white}
}

\newtcolorbox{softbox}[3]{
  enhanced,colback=#2,colframe=#1,borderline west={4pt}{0pt}{#1},
  arc=2mm,boxrule=.4pt,left=4mm,right=4mm,top=3mm,bottom=3mm,
  before skip=6pt,after skip=8pt,title={#3},coltitle=white,
  colbacktitle=#1,fonttitle=\sffamily\bfseries,fontupper=\RaggedRight,
  attach boxed title to top left={xshift=3mm,yshift=-2mm},
  boxed title style={arc=2mm, boxrule=0pt}
}

\newtcolorbox{drawbox}[2]{
  enhanced,colback=white,colframe=#1,arc=4mm,boxrule=1pt,
  left=5mm,right=5mm,top=4mm,bottom=4mm,before skip=8pt,after skip=8pt,
  title={#2},coltitle=white,colbacktitle=#1,fonttitle=\sffamily\bfseries,
  fontupper=\RaggedRight,
  attach boxed title to top left={xshift=4mm,yshift=-2mm},
  boxed title style={arc=3mm, boxrule=0pt}
}

\newtcolorbox{infoband}[1]{
  enhanced,colback=#1!8,colframe=#1!50,arc=2mm,boxrule=.5pt,
  left=4mm,right=4mm,top=2mm,bottom=2mm,before skip=4pt,after skip=4pt,
  fontupper=\small\RaggedRight
}

\newcommand{\linead}[1]{\noindent\color{milcLinea}\rule{#1}{0.5pt}\color{milcNegro}}
\newcommand{\respuesta}[1]{\par\vspace{2mm}\foreach \n in {1,...,#1}{\noindent\color{milcLinea}\rule{\linewidth}{0.5pt}\par\vspace{5mm}}\color{milcNegro}}
\newcommand{\checkbox}{\textcolor{milcGradeDark}{\fbox{\phantom{x}}}\hspace{5pt}}

\newcommand{\phasecard}[4]{
\begin{tcolorbox}[enhanced,colback=#3,colframe=#2,arc=3mm,boxrule=.5pt,height=3.05cm,valign=center,halign=center,left=1.5mm,right=1.5mm,top=2mm,bottom=2mm]
{\sffamily\bfseries\fontsize{22}{24}\selectfont\textcolor{#2}{#1}}\par
{\sffamily\bfseries\small #4}
\end{tcolorbox}
}

\begin{document}
\thispagestyle{empty}
\begin{tikzpicture}[remember picture,overlay]
  \fill[white] (current page.south west) rectangle (current page.north east);
  \path[fill=milcGradePrimary,rounded corners=16mm]
    ([xshift=.55cm,yshift=-.55cm]current page.north west) rectangle
    ([xshift=-.55cm,yshift=3.65cm]current page.south east);

  \node[anchor=north west,text=milcGradeText] at ([xshift=2.0cm,yshift=-1.85cm]current page.north west)
    {\sffamily\bfseries\fontsize{14}{16}\selectfont GRADO};
  \node[anchor=north west,text=milcGradeText,opacity=.82] at ([xshift=2.0cm,yshift=-2.75cm]current page.north west)
    {\sffamily\bfseries\fontsize{9}{11}\selectfont SERIE GUÍAS · TECNOLOGÍA E INFORMÁTICA};

  \begin{scope}[shift={([xshift=-3.05cm,yshift=-2.55cm]current page.north east)}]
    \fill[white,opacity=.80] (-.62,-.52) rectangle (.62,.52);
    \draw[milcGradeText,line width=1pt] (-.62,-.52) rectangle (.62,.52);
    \fill[milcGradeDark] (-.42,-.38) rectangle (-.22,.06);
    \fill[milcGradeAccent] (-.10,-.38) rectangle (.10,.24);
    \fill[milcGradePrimary!60!black] (.22,-.38) rectangle (.42,.38);
  \end{scope}

  \node[anchor=west,text=milcGradeText] at ([xshift=1.9cm,yshift=-12.85cm]current page.north west)
    {\sffamily\bfseries\fontsize{124}{124}\selectfont 8};
  \draw[milcGradeText,line width=1.7pt]
    ([xshift=4.52cm,yshift=-11.68cm]current page.north west) circle (.13cm);

  \node[anchor=west,text=milcGradeText,text width=8.7cm,align=left] at ([xshift=10.35cm,yshift=-11.6cm]current page.north west)
    {
     {\sffamily\bfseries\fontsize{19}{22}\selectfont Lógica avanzada I: estructuras condicionales anidadas (if-else-if) y operadores lógicos AND/OR/NOT}\\[4mm]
     {\sffamily\bfseries\fontsize{23}{26}\selectfont Guía 11-8-TIC}\\[2mm]
     {\sffamily\fontsize{13}{16}\selectfont Periodo Segundo}\\[5mm]
     {\sffamily\bfseries\fontsize{11}{13}\selectfont Institución Educativa Sor María Juliana}\\[1mm]
     {\sffamily\fontsize{10}{12}\selectfont Autor: Dr. Álvaro Cárdenas Orozco}
    };

  \path[fill=milcGradeSoft,rounded corners=7mm]
    ([xshift=1.35cm,yshift=.85cm]current page.south west) rectangle
    ([xshift=-1.35cm,yshift=3.15cm]current page.south east);
  \draw[milcGradeDark,line width=.65pt,rounded corners=7mm]
    ([xshift=1.35cm,yshift=.85cm]current page.south west) rectangle
    ([xshift=-1.35cm,yshift=3.15cm]current page.south east);
  \node[anchor=center,text=milcNegro,text width=17.0cm,align=center] at ([yshift=2.0cm]current page.south)
    {\sffamily\fontsize{10}{12}\selectfont Metodología MILC · Apertura ancestral · Pensamiento computacional · Triángulo Dussel-Estoicismo-Floridi\\
    \textcolor{milcGradeDark}{\bfseries Producto:} Diagrama de flujo + bloque MakeCode con 3 condiciones anidadas + tabla de verdad de 1 expresión compuesta verificada.};
\end{tikzpicture}
\mbox{}
\newpage

\section*{Apertura: Lógica avanzada I: estructuras condicionales anidadas (if-else-if) y operadores lógicos AND/OR/NOT}

\begin{softbox}{milcGradeDark}{milcGradeSoft}{Saber ancestral}
Las \textbf{pruebas del oficio del herrero} en los talleres tradicionales: ¿el hierro está caliente Y bien limpio Y el martillo a la altura correcta? Si falla cualquiera, la pieza se rompe. La decisión depende de varias condiciones que se cruzan al mismo tiempo. Las condiciones lógicas son oficio antes de ser código.
\end{softbox}

\begin{softbox}{milcTurquesa}{milcGris}{Saber contemporáneo}
Las \textbf{estructuras condicionales} en programación (if-else-if) y los operadores lógicos (AND, OR, NOT) traducen al código las decisiones del oficio: \textit{``si esto Y aquello, entonces actúa''}. Son la columna vertebral de cualquier programa que toma decisiones.
\end{softbox}

\begin{softbox}{milcMostaza}{milcGris}{Pregunta-puente}
\textit{¿Qué del juicio del herrero —su capacidad de cruzar varias condiciones antes de actuar— podemos llevar al código que decide?}
\end{softbox}

\begin{softbox}{milcVerde}{milcGris}{Saber y saber hacer}
Al terminar podrás \textbf{escribir condicionales anidadas} (if-else-if), \textbf{combinar} con operadores lógicos AND/OR/NOT, \textbf{construir tablas de verdad} para verificar tu lógica y \textbf{depurar} cuando una condición no se cumple como esperabas.
\end{softbox}

\small
\begin{tabularx}{\textwidth}{>{\bfseries}p{3.0cm}Y}
\rowcolor{milcGradePrimary!12}Autor & Dr. Álvaro Cárdenas Orozco\\
DBA & Aplica estructuras condicionales y operadores lógicos para resolver problemas que requieren múltiples decisiones simultáneas.\\
Referentes & MEN Tecnología · ICFES Pensamiento Computacional · Modelo MILC · CSTA Standards · Floridi (decisiones informacionales) · Estoicismo (lo que depende de mí)\\
Producto final & Diagrama de flujo + bloque MakeCode con 3 condiciones anidadas + tabla de verdad de 1 expresión compuesta verificada.\\
\end{tabularx}
\normalsize

\begin{titlebox}{milcGradeDark}
Ruta MILC: así vamos a aprender
\end{titlebox}
\begin{tabularx}{\textwidth}{>{\centering\arraybackslash}X>{\centering\arraybackslash}X>{\centering\arraybackslash}X>{\centering\arraybackslash}X}
\phasecard{1}{milcVerde}{milcGris}{Escuta\\[-1mm]\small Observo, escucho y pregunto.} &
\phasecard{2}{milcTurquesa}{milcGris}{Sistematización\\[-1mm]\small Pensamiento computacional.} &
\phasecard{3}{milcMagenta}{milcGris}{Praxis\\[-1mm]\small Construyo y aplico.} &
\phasecard{4}{milcVino}{milcGris}{Evaluación\\[-1mm]\small Reflexión liberadora.}
\end{tabularx}


\begin{titlebox}{milcVerde}
Fase 1 · Escuta
\end{titlebox}

\begin{softbox}{milcVerde}{milcGris}{Escena para escuchar}
Felipe modela en MakeCode el semáforo del cruce frente al colegio. La regla es: SI hay peatón Y la luz está en verde para autos, los autos paran. Pero se da cuenta que falta otra condición: ¿qué pasa si hay peatón pero está en la acera, no cruzando? La condición simple no alcanza. Necesita anidar y combinar lógica.
\end{softbox}

\checkbox Recuerdo una decisión real donde tuve que cruzar varias condiciones a la vez.\\
\checkbox Identifico la diferencia entre AND (Y) y OR (O) en español cotidiano.\\
\checkbox Distingo cuándo conviene anidar condicionales (if-else-if) en lugar de varias separadas.

\textbf{Una decisión de mi vida con varias condiciones a la vez fue:}\\[1mm]
\linead{\linewidth}\\[1mm]
\linead{\linewidth}

\textbf{Las condiciones que tuve que cruzar fueron:}\\[1mm]
\linead{\linewidth}\\[1mm]
\linead{\linewidth}

\textbf{Si la programara, mi pseudocódigo empezaría con:}\\[1mm]
\linead{\linewidth}\\[1mm]
\linead{\linewidth}

\begin{infoband}{milcVerde}
\textbf{Saber más.} Los operadores lógicos vienen de George Boole (s. XIX). Su "álgebra booleana" es la base de TODA la computación moderna: cada decisión que toma tu celular se reduce a combinaciones de AND, OR, NOT. Lo que el herrero hacía con intuición, hoy se calcula con lógica formal.
\end{infoband}

\begin{titlebox}{milcTurquesa}
Fase 2 · Sistematización con pensamiento computacional
\end{titlebox}

\begin{softbox}{milcTurquesa}{milcGris}{Cuatro pilares aplicados al tema}
El pensamiento computacional descompone la decisión en sus condiciones, reconoce patrones de combinación lógica, abstrae la regla \textit{toda decisión es función de variables} y diseña un algoritmo verificable con tabla de verdad.
\end{softbox}

\small
\begin{tabularx}{\textwidth}{>{\bfseries\RaggedRight\arraybackslash}p{3.6cm}Y}
\rowcolor{milcTurquesa!90}\color{white}Pilar & \color{white}\bfseries Aplicación al tema\\
1. Descomposición & Descompón la decisión: ¿cuántas condiciones hay? ¿se necesitan todas (AND) o basta una (OR)? ¿hay condiciones que invierten (NOT)?\\
2. Reconocimiento de patrones & Patrones: AND para cruce obligatorio (todas verdaderas). OR para alternativa (al menos una). NOT para inversión.\\
3. Abstracción & Abstracción: la condición compuesta se puede escribir como UNA expresión booleana que evalúa V o F.\\
4. Algoritmo conceptual & Algoritmo: \textbf{1)} listar condiciones · \textbf{2)} elegir operadores · \textbf{3)} construir expresión · \textbf{4)} crear tabla de verdad · \textbf{5)} verificar todos los casos · \textbf{6)} traducir a código.\\
\end{tabularx}
\normalsize

\begin{drawbox}{milcTurquesa}{Anatomía de una expresión booleana}
\textbf{Variables (proposiciones):} cada condición que evalúa V/F. Ej: \texttt{hay\_peaton}, \texttt{luz\_verde}. \\
\textbf{Operador AND (Y):} ambas verdaderas. Ej: \texttt{hay\_peaton AND luz\_verde}. \\
\textbf{Operador OR (O):} al menos una verdadera. Ej: \texttt{lluvia OR frio}. \\
\textbf{Operador NOT (no):} invierte. Ej: \texttt{NOT hay\_peaton}. \\
\textbf{Paréntesis:} agrupan y fuerzan orden. Ej: \texttt{(A OR B) AND C}. \\
\textbf{Tabla de verdad:} todos los casos posibles para verificar la expresión.
\end{drawbox}

\begin{softbox}{milcMostaza}{milcGris}{Errores comunes y su corrección}
\begin{itemize}
\item Confundir AND con OR en lenguaje natural $\rightarrow$ ``si llueve y hace frío'' (AND) vs ``si llueve o hace frío'' (OR) son DOS reglas distintas.
\item Olvidar paréntesis $\rightarrow$ \texttt{A OR B AND C} se evalúa diferente a \texttt{(A OR B) AND C}.
\item No probar todos los casos $\rightarrow$ una tabla de verdad de 3 variables tiene 8 filas; verifícalas todas.
\item Anidar demasiado $\rightarrow$ if dentro de if dentro de if se vuelve ilegible; usa AND/OR cuando puedas.
\item Negar mal $\rightarrow$ \texttt{NOT (A AND B)} es \texttt{NOT A OR NOT B}, no \texttt{NOT A AND NOT B}.
\end{itemize}
\end{softbox}

\textbf{Mi expresión booleana del semáforo del cruce sería:}\\[1mm]
\linead{\linewidth}\\[1mm]
\linead{\linewidth}

\textbf{Construida la tabla de verdad, los casos verdaderos son:}\\[1mm]
\linead{\linewidth}\\[1mm]
\linead{\linewidth}

\begin{titlebox}{milcMagenta}
Fase 3 · Praxis con pensamiento computacional
\end{titlebox}

\begin{softbox}{milcMagenta}{milcGris}{Construyendo el producto con los cuatro pilares}
Construir un sistema de decisión con condicionales anidadas es un sistema. Decomponemos las condiciones, reconocemos patrones lógicos, abstraemos la expresión y seguimos un algoritmo verificable.
\end{softbox}

\small
\begin{tabularx}{\textwidth}{>{\bfseries\RaggedRight\arraybackslash}p{3.6cm}Y}
\rowcolor{milcMagenta!90}\color{white}Pilar & \color{white}\bfseries Aplicación al producto\\
Descomposición & Sistema: 1) variables · 2) condiciones · 3) operadores · 4) anidamiento · 5) tabla de verdad · 6) código.\\
Patrones & Patrones de buenas decisiones programadas: condición clara, anidamiento mínimo, tabla de verdad documentada.\\
Abstracción & Función esencial: el código toma UNA decisión correctamente para todos los casos posibles.\\
Algoritmo de armado & Pasos: definir variables $\rightarrow$ escribir condiciones $\rightarrow$ combinar con operadores $\rightarrow$ verificar con tabla $\rightarrow$ traducir a MakeCode.\\
\end{tabularx}
\normalsize

\begin{drawbox}{milcMagenta}{Checklist del sistema de decisión}
\checkbox 3 variables booleanas identificadas claramente.\\
\checkbox Expresión compuesta usando AND, OR y NOT.\\
\checkbox Tabla de verdad completa (2\^{}n filas).\\
\checkbox Diagrama de flujo del algoritmo.\\
\checkbox Bloque MakeCode con la lógica.\\
\checkbox Prueba con al menos 3 casos distintos.
\end{drawbox}

\begin{softbox}{milcMagenta}{milcGris}{Plantilla de trabajo}
\textbf{Variables:} A = \_\_\_, B = \_\_\_, C = \_\_\_. \\
\textbf{Expresión:} \texttt{(A AND B) OR (NOT C)} (ejemplo). \\
\textbf{Tabla de verdad:} 8 filas (2\^{}3 = 8 casos posibles). \\
\textbf{MakeCode:} \texttt{if (A and B) or (not C)} then \textit{acción}. \\
\textbf{Prueba:} caso 1 (A=V, B=V, C=F) $\rightarrow$ V. caso 2 (A=F, B=F, C=V) $\rightarrow$ F. \\
\textbf{Documentación:} explicar qué representa cada variable en el contexto real.
\end{softbox}

\textbf{Mi sistema de decisión: variables + expresión + tabla + código MakeCode:}\\[1mm]
\linead{\linewidth}\\[1mm]
\linead{\linewidth}\\[1mm]
\linead{\linewidth}

\begin{titlebox}{milcMostaza}
Producto de la sesión
\end{titlebox}
\textbf{Diagrama de flujo + bloque MakeCode + tabla de verdad de expresión compuesta}.
\begin{softbox}{milcMostaza}{milcGris}{Criterios de calidad}
\begin{itemize}
\item 3 variables booleanas con significado real.
\item Expresión usando AND, OR, NOT con paréntesis correctos.
\item Tabla de verdad completa y verificada.
\item Diagrama de flujo legible.
\item Bloque MakeCode funcional.
\item Documentación de qué hace cada parte.
\end{itemize}
\end{softbox}

\begin{titlebox}{milcVino}
Fase 4 · Evaluación liberadora — cinco dimensiones
\end{titlebox}

\begin{softbox}{milcVino}{milcGris}{Sentido de la evaluación}
Evaluar no es solo revisar una nota. Es reconocer qué aprendí, qué cambió en mí, qué emociones manejé, cómo me relaciono con la ciudad y la comunidad, y qué le dejaré a quien viene detrás.
\end{softbox}

\textbf{Mi reflexión liberadora en cinco dimensiones.} Completa cada frase iniciada:

\small
\begin{tabularx}{\textwidth}{>{\bfseries\RaggedRight\arraybackslash}p{3.4cm}Y>{\RaggedRight\arraybackslash}p{4.6cm}}
\rowcolor{milcVino}\color{white}Dimensión & \color{white}\bfseries Mi respuesta (frame starter) & \color{white}\bfseries Algo que descubrí\\
1. Personal & Lo que más cambió en mí fue \linead{6.5cm} & \linead{4.2cm}\\
2. Emocional & La emoción más fuerte fue \linead{2.5cm} y la regulé \linead{3cm} & \linead{4.2cm}\\
3. Ciudadana & En democracia esto importa porque \linead{6cm} & \linead{4.2cm}\\
4. Local (Cartago) & En mi barrio sería útil para \linead{6.5cm} & \linead{4.2cm}\\
5. Intergeneracional & A mi abuela/abuelo le diría \linead{6.5cm} & \linead{4.2cm}\\
\end{tabularx}
\normalsize

\begin{infoband}{milcVino}
\textbf{Para la próxima sustentación.} Vuelve a esta tabla en seis meses. Verás que las respuestas cambian — no porque lo aprendido cambie, sino porque tú cambias. La evaluación liberadora es una conversación contigo a través del tiempo.
\end{infoband}


\begin{titlebox}{milcGradeDark}
Triángulo de pensamiento · cierre filosófico
\end{titlebox}

\begin{softbox}{milcVino}{milcGris}{Dussel · lente del nosotros}
\textit{``La analéctica supera a la dialéctica al partir del Otro como Otro.''}

Toda regla lógica excluye casos que no contempla. Una condición \texttt{if mayor\_de\_edad} excluye al adolescente que necesita una decisión urgente. La analéctica con código pregunta: ¿qué caso real quedó fuera de mi expresión booleana?

\textbf{¿Qué caso real (persona, situación) NO captura tu expresión condicional? ¿Cómo lo incluirías?}
\end{softbox}
\linead{\linewidth}\\[1mm]
\linead{\linewidth}

\begin{softbox}{milcMostaza}{milcGris}{Estoicismo · Marco Aurelio · cuidado interior}
\textit{``El obstáculo en el camino se vuelve el camino.''}

Cada error de lógica (un AND donde debías poner OR) es maestro: la tabla de verdad te muestra dónde fallaste y por qué. La depuración es ejercicio estoico de aprender del error.

\textbf{¿Qué error de lógica reciente se convirtió en aprendizaje permanente?}
\end{softbox}
\linead{\linewidth}\\[1mm]
\linead{\linewidth}

\begin{softbox}{milcTurquesa}{milcGris}{Floridi · lente de la infoesfera}
\textit{``Cada decisión de información tiene peso ético.''}

Una condición mal escrita en un software de salud (ej: \texttt{if dosis > 100} en lugar de \texttt{if dosis < 100}) puede dañar a alguien real. La lógica no es ejercicio académico — es responsabilidad sobre vidas.

\textbf{Si tu sistema de decisión se aplicara a una situación real (semáforo, alarma médica), ¿quién se afectaría si fallas la condición?}
\end{softbox}
\linead{\linewidth}\\[1mm]
\linead{\linewidth}

\begin{titlebox}{milcVino}
Compromiso final
\end{titlebox}

\small
\begin{tabularx}{\textwidth}{>{\RaggedRight\arraybackslash}p{3.0cm}YYY}
\rowcolor{milcVino}\color{white}\bfseries Criterio & \color{white}\bfseries Lo logré & \color{white}\bfseries En proceso & \color{white}\bfseries Necesito apoyo\\
Comprensión del tema & Explico ideas centrales con mis palabras. & Reconozco ideas pero falta precisión. & Necesito repasar conceptos base.\\
Pensamiento computacional & Apliqué los 4 pilares al diseñar mi producto. & Apliqué algunos pilares. & Necesito releer la fase 2.\\
Aplicación y producto & Mi producto cumple los criterios y se puede revisar. & Producto funcional con ajustes pendientes. & Aún en construcción.\\
Reflexión 5 dimensiones & Respondí cada dimensión con honestidad. & Respondí algunas con detalle. & Mis respuestas son muy generales.\\
Triángulo de pensamiento & Conecté las 3 voces con mi trabajo real. & Conecté al menos una. & No relaciono las voces con mi proceso.\\
\end{tabularx}
\normalsize

\textbf{Hoy entendí que:}\\[1mm]
\linead{\linewidth}\\[1mm]
\linead{\linewidth}

\textbf{Mi compromiso para la próxima sesión es:}\\[1mm]
\linead{\linewidth}\\[1mm]
\linead{\linewidth}

\begin{softbox}{milcMostaza}{milcGris}{Cita de despedida · Marco Aurelio}
\textit{``El obstáculo en el camino se vuelve el camino. Lo que estorba al camino se vuelve camino.''}\\[1mm]
Cada error de hoy, bien observado, se convierte en pieza del camino que pisarás mañana.
\end{softbox}

\small
\begin{tabularx}{\textwidth}{>{\bfseries}p{3.6cm}Y}
\rowcolor{milcGradeDark!12}Nombre completo: & \linead{10cm}\\
Fecha: & \linead{4cm} \quad Curso/grupo: \linead{4cm}\\
Mi firma: & \linead{9cm}\\
\end{tabularx}
\normalsize

\begin{infoband}{milcGradeDark}
\textbf{Próxima sesión.} Profundizaremos en tablas de verdad y condiciones compuestas. La lógica del herrero se vuelve disciplina booleana.
\end{infoband}

\end{document}
